\section{产品风险：低伤世界也可能变成高伤世界}

Elys 的愿景是低伤世界：更少摩擦、更妙连接。但任何社交产品都可能制造伤害，AI 社交还会放大这种风险。分身可能误代表用户，context 可能泄露，AI 可能制造虚假热闹，用户可能被算法推向过度暴露或不健康关系。

\subsection{风险清单}

\noindent\begin{tabularx}{\textwidth}{p{0.22\textwidth}p{0.34\textwidth}X}
\toprule
风险 & 表现 & 需要的机制 \\
\midrule
误代表 & 分身说出用户不认可的话 & 预览、撤回、敏感动作确认。 \\
隐私泄露 & 高维 context 被平台或他人滥用 & 本地存储、加密、分级授权。 \\
虚假社交 & AI 互动淹没真人互动 & 标记 AI/真人，优化真人连接率。 \\
操纵与成瘾 & 分身为了留存制造情绪依赖 & 透明目标、健康使用限制。 \\
关系误配 & 高维匹配带来错误亲密感 & 逐步确认、真人反馈、边界提示。 \\
\bottomrule
\end{tabularx}

\begin{warningbox}{低伤世界不是自动结果}
AI 降低摩擦，也可能降低防线。产品必须把权限、透明度、真人确认和健康边界设计进去，否则低伤世界可能变成更隐蔽的高伤世界。
\end{warningbox}

\subsection{本章小结}

AI 社交的风险和价值来自同一件事：它更懂你，也更主动。理解这一点，才能设计出既有连接能力又有边界感的产品。


